
\documentclass[10pt,a4paper]{article}
\usepackage[margin=0.65in]{geometry}
\usepackage[T1]{fontenc}
\usepackage[utf8]{inputenc}
\usepackage{lmodern}
\usepackage{microtype}
\usepackage{multicol}
\usepackage{booktabs}
\usepackage{array}
\usepackage{amsmath}
\usepackage[table]{xcolor}
\usepackage{graphicx}
\usepackage{tikz}
\usepackage{hyperref}
\setlength{\parindent}{0pt}
\setlength{\parskip}{0.35em}
\pagestyle{empty}
\definecolor{accent}{named}{black}
\begin{document}
\begin{center}
{\LARGE\bfseries MDPI Family}\par
\vspace{0.25em}
{\small Journal Manuscript Preview Library}\par
\vspace{0.15em}
{\footnotesize Family baseline preview for MDPI Family}\par
\vspace{0.2em}
{\footnotesize Profile slug: mdpi \quad Family: mdpi \quad Scope: Family baseline}
\end{center}
\vspace{0.2em}
\noindent{\color{red!65!black}\bfseries Article Type:} Placeholder Research Article\par\vspace{0.4em}


\noindent\textbf{Authors.} A. Placeholder, B. Example, and C. Template\par
\noindent\textbf{Affiliations.} Department of Simulated Formatting, Placeholder Institute\par
\noindent\textbf{Contact.} preview@example.org\par
\vspace{0.4em}
\fcolorbox{black!45}{gray!8}{\begin{minipage}{0.96\linewidth}\textbf{Abstract.} This simulated first-page preview uses placeholder text, a placeholder figure, and a placeholder table to visualize how a manuscript may look under the selected journal profile. The content is intentionally generic so attention stays on spacing, title treatment, abstract presentation, caption location, and the overall page rhythm.\end{minipage}}\par\vspace{0.6em}
\noindent\textbf{Keywords:} placeholder manuscript, journal layout, figure slot, table slot, references.\par\vspace{0.6em}


\small
\section*{1. Introduction}
This placeholder article body is intentionally structured as a compact submission-style manuscript so the rendered page is easier to inspect as a serious journal preview. The opening section focuses on how the title block, abstract, keywords, and first paragraphs combine into an overall page rhythm that resembles a practical submission rather than a poster-like mockup.

\section*{2. Related Work}
The related-work section is included even in starter previews because it makes the hierarchy more realistic and helps the reader judge whether literature framing, method description, and experimental content are clearly separated.

\section*{3. Methodology}
The methodology section introduces representative technical material so the preview can expose equation spacing, paragraph density, caption rhythm, and the visual balance of the first major figure.

\begin{equation}
\mathcal{L} = \sum_{i=1}^{M} w_i \left\| y_i - \hat{y}_i \right\|_2^2,
\end{equation}
where $y_i$ denotes a placeholder reference output, $\hat{y}_i$ denotes the estimated output, and $w_i$ denotes a weighting coefficient included only to simulate realistic notation.


\begin{center}\fcolorbox{black!60}{gray!10}{\begin{minipage}[c][2.25cm][c]{0.93\linewidth}\centering\Large Image Placeholder\\[0.4em]\normalsize Simulated figure area\end{minipage}}\end{center}\noindent\textbf{Figure 1.} Placeholder figure caption showing the expected caption style and relative spacing.\par\vspace{0.5em}

\section*{4. Experiments}
The experiments section is expanded so the preview can show a more realistic transition from methodology to quantitative evidence. This makes the resulting PNG materially easier to inspect for spacing, float handling, and table readability.
\textbf{Experimental setup.} This paragraph stands in for datasets, hardware conditions, baseline models, and protocol notes that are typically required before the main comparison table appears.

\noindent\textbf{Table 1.} Placeholder table caption showing metric labels and compact journal spacing.\par\begin{center}
\begin{tabular}{lccc}
\toprule
Method & Score & Error & Time\\
\midrule
Placeholder A & 0.91 & 0.08 & 1.4\\
Placeholder B & 0.88 & 0.11 & 1.7\\
Placeholder C & 0.84 & 0.15 & 2.0\\
\bottomrule
\end{tabular}\end{center}\vspace{0.5em}

\textbf{Quantitative discussion.} The surrounding text exists to make the table feel integrated into a real manuscript page instead of floating as an isolated placeholder. It also lets the preview expose line lengths, caption spacing, and how the figure-table pair reads at a glance.

\section*{5. Conclusion and Discussion}
The final section is intentionally brief but still substantial enough to reveal whether the page closes cleanly after the main figure and table. The content remains generic by design, while the structure aims to make cross-journal layout comparison faster and more trustworthy.


\noindent \textbf{Data Availability.} Placeholder data and code would be linked here.\par
\noindent \textbf{Author Contributions.} Conceptualization, analysis, writing, and review are listed here as placeholders.\par
\noindent \textbf{Competing Interests.} The authors declare no competing interests in this simulated preview.\par
\section*{6. References}
\vspace{0.5em}
\noindent \textbf{References.} [1] A. Smith and B. Lee, ``Journal template adaptation,'' \textit{Placeholder Journal}, 2024. [2] C. Garcia \textit{et al.}, ``Figure and table policy notes,'' 2025.\par
\vspace{0.6em}
\fcolorbox{black!45}{gray!6}{\begin{minipage}{0.98\linewidth}
\footnotesize Simulated preview built from the local starter profile library.\\ Use the official template to finalize submission-ready formatting.
\end{minipage}}
\end{document}
